% Reader-first, claim-preserving front matter for the Problem 593 manuscript.
% This fragment is not included automatically.  It is written for the
% publication branch containing the finite parameter spectrum.

\begin{abstract}
A finite triple system is \emph{obligatory} if it embeds in every triple system
of uncountable chromatic number.  We give an alternative proof of the complete
classification.  After isolated vertices are removed, obligatoriness is
equivalent to the following three finite conditions: the system is linear,
every hyperedge-node of its Levi graph is incident with a bridge, and every
Berge cycle has even length.  Equivalently, the obligatory systems are obtained
from private-vertex expansions of finite bipartite graphs and finite edgeless
systems by finite disjoint unions and one-point amalgamations.  The proof
separates the finite bridge-block decomposition from the infinitary host
constructions: deleting the Levi bridges produces bipartite expansion pieces
attached along a quotient forest, while a one-apex lift supplies avoiding hosts
for the missing-bridge and odd-cycle obstructions.  The same decomposition gives
sharp finite parameter consequences.  A reduced obligatory system with
$m\ge1$ hyperedges and $1\le c\le m$ connected components has order $n$ if and
only if
\[
  m+2(c-1)+\left\lceil2\sqrt{m-c+1}\right\rceil
  \le n\le 2m+c.
\]
A Lean~4 development verifies both directions of the finite classification and
the isolated-vertex reduction.  All arguments are in ZFC.
\end{abstract}

\section*{Introduction}\label{introduction-reader-first}

The hosts in Erd\H{o}s Problem~593 are infinite, but the answer is governed by a
finite incidence graph.  A finite triple system $F$ is \emph{obligatory} if
every triple system of uncountable chromatic number contains an injective,
not-necessarily-induced copy of $F$.  Erd\H{o}s asked for a characterisation of
these systems~\citep{erdos1995}; the problem is catalogued as Erd\H{o}s
Problem~593~\citep{bloom593}.

For a finite graph $J$, let $J^+$ be its private-vertex expansion: each graph
edge $uv$ is replaced by the triple $\{u,v,p_{uv}\}$, where the points $p_{uv}$
are new and pairwise distinct.  Let $\mathcal B$ be the smallest class
containing $J^+$ for every finite bipartite graph $J$, together with every
finite edgeless triple system, and closed under finite disjoint unions and
one-point amalgamations.  For a triple system $F$, write $F^\circ$ for the
system obtained by deleting its isolated vertices, and write $I(F)$ for its
Levi graph.

\begin{theoremA}\label{theorem-a-reader-first}
For every finite triple system $F$, the following are equivalent.
\begin{enumerate}
\item $F$ is obligatory;
\item $F^\circ\in\mathcal B$;
\item $F^\circ$ is linear, every hyperedge-node of $I(F^\circ)$ is incident
with a bridge, and every Berge cycle of $F^\circ$ has even length.
\end{enumerate}
\end{theoremA}

The proof is best understood as three independent implications:
\[
  \mathcal B\Longrightarrow\text{obligatory},
  \qquad
  \mathcal B\Longleftrightarrow\text{intrinsic},
  \qquad
  \neg\text{intrinsic}\Longrightarrow\neg\text{obligatory}.
\]
The middle equivalence is the finite geometric core.  Delete every bridge from
the Levi graph.  A hyperedge-node has degree three and, by hypothesis, loses at
least one incidence; its remaining degree is therefore zero or two.  Each
active bridge-free component suppresses to a simple graph.  Linearity rules out
parallel edges, and evenness of Berge cycles makes the graph bipartite.  The
removed incidences supply the private points, so each component is an expansion
$J^+$.  Contracting the bridge-free components turns the original bridges into
a forest, whose rooted order reconstructs $F^\circ$ by disjoint unions and
one-point amalgamations.

The positive implication shows that these pieces and attachment operations are
forced in every uncountably chromatic host.  The expansion theorem of
Reiher~\citep[Theorem~1.2]{reiher2024} supplies the basic atoms; we also include
a direct proof in the notation of this paper.  Closure under disjoint union is
immediate after finite deletion.  Closure under one-point amalgamation follows
from a rooted-abundance lemma: in an uncountably chromatic host, the set of
vertices supporting arbitrarily many almost-disjoint rooted copies of an
obligatory system remains uncountably chromatic.

For necessity, the intrinsic test can fail in exactly three ways.  A classical
Erd\H{o}s--Rado construction gives an uncountably chromatic linear triple system
and therefore excludes every nonlinear $F$.  The remaining two obstructions
are controlled by the complete-rank one-apex lift.  A finite linear trace in the
lift decomposes into base fibres, each isomorphic to an expansion $J_s^+$; two
fibres meet in at most one point, and their support-incidence graph is a forest.
Consequently every hyperedge-node in the trace is incident with a bridge, and
every Berge cycle projects with the same length to the base graph.  Taking the
base graph to be $K_{\omega_1}$ excludes a missing bridge, while an
uncountably chromatic graph of sufficiently large odd girth excludes an odd
Berge cycle.

Li's preprint contains the first publicly posted complete proof of
Theorem~\ref{theorem-a-reader-first} and introduces the complete-rank one-apex
lift and bridge-trace method used in the negative direction
\citep[Theorem~1.1, Sections~3--4]{li2026}.  The present paper gives an
alternative implementation, a direct all-bridges decomposition of the finite
structure, sharp finite parameter consequences, and a Lean formalisation of
this implementation.  No competing priority claim is made for the
classification or the shared lift architecture.

The bridge-block argument retains numerical information.  It produces an
ordinary bipartite shadow $J$ with the same number of edges and connected
components as $F^\circ$ and with
\[
  |V(J)|=|V(F^\circ)|-|E(F^\circ)|.
\]
Combining this identity with the exact order--size range for finite bipartite
graphs gives the following sharp spectrum.

\begin{theorem}[Order--size--component spectrum]\label{theorem-spectrum-reader-first}
Let $m\ge1$ and $1\le c\le m$.  A reduced obligatory triple system with
$m$ hyperedges, $n$ vertices, and exactly $c$ connected components exists if
and only if
\[
  m+2(c-1)+\left\lceil2\sqrt{m-c+1}\right\rceil
  \le n\le 2m+c.
\]
Every integer in this interval occurs.
\end{theorem}

The manuscript is organised to keep the finite geometry visible.  After the
preliminaries, we prove the bridge-block equivalence between the constructive
and intrinsic descriptions.  We then establish the positive direction, define
the one-apex lift, and prove the finite trace theorem used by the three avoiding
hosts.  The final mathematical section derives the bipartite shadow and its
finite parameter consequences.  The Lean scope, reproducibility information,
and provenance statements are collected at the end rather than repeated inside
the proofs.
